\subsection{由完备赋范李代数定义的群}

设 \(H = \sum_{r,s \geq 0} H_{r,s} \in \hat{L}(\{U, V\})\) 为豪斯多夫级数（§ 6，第 4 号，定义 1）。我们将证明相应的带连续分量的形式幂级数

\[
\tilde {\mathrm{H}} = \sum_ {r, s \geqslant 0} \tilde {\mathrm{H}} _ {r, s} \in \hat {\mathrm{P}} (\mathfrak {g} \times \mathfrak {g}, \mathfrak {g})\tag{10}
\]

是收敛的（《微分流形与解析流形》R，4.1.1）。

令 \(r \geqslant 0\)、\(s \geqslant 0\) 为满足 \(r + s \neq 0\) 的数，并令 \(\| \tilde{\mathrm{H}}_{r,s} \|\) 为连续多项式 \(\tilde{\mathrm{H}}_{r,s}\) 的范数（《微分流形与解析流形》R 附录第 2 号）。

引理 3。

\[
\left\| \tilde {\mathrm{H}} _ {r, s} \right\| \leqslant a ^ {- (r + s - 1) \theta}.
\]

设 B 是关于 I 的 Hall 集，令 \(H_{r,s} = \sum_{b \in B} \lambda_b e_b\) 为 \(H_{r,s}\) 关于 \(L(\{U, V\})\) 中相应基的分解。则

\[
\left| \lambda_ {b} \right| \leqslant a ^ {- (r + s - 1) \theta}.\tag{11}
\]

当 \(p = 0\) 时这显然成立，因为 \(\lambda_b \in \mathbf{Q}\)；当 \(p \neq 0\) 时则由第 1 号命题 1 得到。

此外，

\[
\left\| \tilde {e} _ {b} \right\| \leqslant 1 \quad \text { for } b \in \mathbf {B}.\tag{12}
\]

更一般地，通过对 \(n\) 作归纳可证明：对于关于两个不定元 U 和 V 的每个交错式 \(b\)，其次数为 \(n\)（§2，第 6 号），都有 \(\| \tilde{b} \| \leqslant 1\)。若 \(n = 1\)，则 \(\tilde{b}\) 是从 \(\mathfrak{g} \times \mathfrak{g}\) 到 \(\mathfrak{g}\) 的投影之一，因而范数 \(\leqslant 1\)；若 \(n > 1\)，则存在两个交错式 \(b_1\) 和 \(b_2\)，它们的次数均为 \(< n\)，使得 \(b = [b_1, b_2]\)。由于映射 \(\gamma: (x, y) \mapsto [x, y]\) 是从 \(\mathfrak{g} \times \mathfrak{g}\) 到 \(\mathfrak{g}\) 的范数 \(\leqslant 1\) 的双线性映射，有（《微分流形与解析流形》R，附录，第 4 号）

\[
\| \tilde {b} \| = \| \gamma \circ (\tilde {b} _ {1}, \tilde {b} _ {2}) \| \leqslant \| \tilde {b} _ {1} \|. \| \tilde {b} _ {2} \| \leqslant 1.\tag{13}
\]

关系 (11) 和 (12) 蕴含引理。

命题 2。形式幂级数 \(\tilde{\mathbf{H}}\) 是收敛级数（《微分流形与解析流形》R，4.1.1）。若 \(\mathbf{G}\) 是球 \(\{x\in\mathfrak{g}\mid\|x\|<a^{\theta}\}\)，则 \(\tilde{\mathbf{H}}\) 的绝对收敛域（《微分流形与解析流形》R，4.1.3）包含 \(\mathbf{G}\times\mathbf{G}\)。

若 \(u\) 和 \(v\) 是满足 \(>0\) 且 \(u < a^{\theta}\) 的两个 \(v < a^{\theta}\) 的实数，则由引理 3

\[
\| \tilde {\mathrm{H}} _ {r, s} \| u ^ {r} v ^ {s} \leqslant a ^ {\theta} (u a ^ {- \theta}) ^ {r} (v a ^ {- \theta}) ^ {s}\tag{14}
\]

并且当 \(\| \tilde{\mathrm{H}}_{r,s}\| u^r v^s\) 趋于无穷大时，\(r + s\) 趋于 0。

记 \(h \colon G \times G \to \mathfrak{g}\) 为由 \(\tilde{\mathbf{H}}\) 定义的解析函数（《微分流形与解析流形》R，4.2.4），即由公式

\[
h (x, y) = \sum_ {r, s \geqslant 0} \tilde {\mathrm{H}} _ {r, s} (x, y) = \sum_ {r, s \geqslant 0} \mathrm{H} _ {r, s} (x, y) \quad \text { for } (x, y) \in \mathrm{G} \times \mathrm{G}.\tag{15}
\]

此函数称为 g 的豪斯多夫函数。

令 \((x,y)\in \mathbf{G}\times \mathbf{G}\)。则

\begin{enumerate}
\def\labelenumi{(\arabic{enumi})}
\setcounter{enumi}{15}
\tightlist
\item
\end{enumerate}

\[
\left\| \tilde {\mathrm{H}} _ {r, s} (x, y) \right\| \leqslant \sup (\left\| x \right\|, \left\| y \right\|)\tag{16}
\]

\[
\| h (x, y) \| \leqslant \sup (\| x \|, \| y \|).\tag{17}
\]

\begin{enumerate}
\def\labelenumi{(\arabic{enumi})}
\setcounter{enumi}{16}
\tightlist
\item
  由 (16) 立即得到，而 \(r = s = 0\) 时 (16) 显然成立；若 \(r \geqslant 1\)，则
\end{enumerate}

\[
\begin{array}{r l} \| \tilde {\mathrm{H}} _ {r, s} (x, y) \| & \leqslant \| \tilde {\mathrm{H}} _ {r, s} \| \| x \| ^ {r} \| y \| ^ {s} \\ & \leqslant \| x \| \left(\frac {\| x \|}{a ^ {\theta}}\right) ^ {r - 1} \left(\frac {\| y \|}{a ^ {\theta}}\right) ^ {s} \\ & \leqslant \| x \|; \end{array}
\]

若 \(s \geqslant 1\)，同理可证。

特别地，对于 \(\| h(x,y)\| < a^{\theta}\)，有 \((x,y)\in \mathbf{G}\times \mathbf{G}\)。

命题 3。令 G 为球 \(\{x\in\mathfrak{g}\mid\|x\|<a^{\theta}\}\)。解析映射

\[
h \colon \mathrm{G} \times \mathrm{G} \rightarrow \mathrm{G}
\]

使 G 成为一个群，其中 0 是单位元，并且对所有 \(x \in G\)，-x 是 x 的逆元。此外，若 R 是满足 \(0 < R < a^{\theta}\) 的实数，则球

\[
\{x \in \mathfrak {g} | \| x \| <   \mathrm{R} \}
\]

（分别为 \(\{x \in \mathfrak{g} \mid \|x\| \leqslant R\}\)）是 G 的开子群。

由于 \(\mathrm{H}(\mathrm{U}, -\mathrm{U}) = 0\) 且 \(\mathrm{H}(0, \mathrm{U}) = \mathrm{H}(\mathrm{U}, 0) = \mathrm{U}\)，有 \(h(x, -x) = 0\)，并且

\[
h (0, x) = h (x, 0) = x
\]

对所有 \(x \in G\) 成立。因此还需证明结合律公式

\[
h (h (x, y), z) = h (x, h (y, z)) \quad \text { for } x, y, z \text { in } G.\tag{18}
\]

由于

\[
\mathrm{H} (\mathrm{H} (\mathrm{U}, \mathrm{V}), \mathrm{W}) = \mathrm{H} (\mathrm{U}, \mathrm{H} (\mathrm{V}, \mathrm{W}))
\]

在 \(\hat{\mathbf{L}} (\{\mathrm{U},\mathrm{V},\mathrm{W}\})\) 中（§ 6，第 5 号，命题 4），有

\[
\tilde {\mathrm{H}} \circ (\tilde {\mathrm{H}} \times \mathrm{Id} _ {\mathfrak {g}}) = \tilde {\mathrm{H}} \circ (\mathrm{Id} _ {\mathfrak {g}} \times \tilde {\mathrm{H}})\tag{19}
\]

在 \(\hat{\mathbf{P}} (\mathfrak{g}\times \mathfrak{g}\times \mathfrak{g};\mathfrak{g})\) 中（第 2 号），并且由 (16) 及《微分流形与解析流形》R，4.1.5，(19) 蕴含 (18)。

*换言之（第 III 章，§1），带有豪斯多夫函数的 G 是一个李群。*

